\subsection{Characteristic finite targets for the three weight-five special directions}

\begin{statement}{D.37}{Proved}{An exact characteristic Magnus target and all three ordinary outer parameters}
For every \(n\geq1\) there is an explicitly specified characteristic
finite quotient
\[
 \rho_n:\widehat F\twoheadrightarrow P_n,\qquad
 |P_n|=5^{17n+2},
 \tag{D.37.1}
\]
which is metabelian of nilpotency class six.  It has
\(P_n^{\rm ab}\simeq(\mathbf Z/5^{n+1})^2\),
derived subgroup \((\mathbf Z/5^n)^{15}\), and centre
\(\gamma_6P_n\simeq(\mathbf Z/5^n)^5\).
Its kernel is fully invariant, by an explicit Fox substitution map.
Coefficient reduction gives onto maps \(P_{n+1}\to P_n\).

There is a specified injective homomorphism
\[
 (\mathbf Z/5^n)^3\longrightarrow\operatorname{Out}^{\sharp}(P_n),
 \qquad (A,B,C)\longmapsto[\beta_{A,B,C}],
 \tag{D.37.2}
\]
whose representatives fix \(x_1\) and have actual multiplier-one
conjugators at \(x_2\) and \(x_3\).
The parameters have exactly the requested leading coordinate
\[
 (X+Y)(AX^2+BXY+CY^2)[x,y].
 \tag{D.37.3}
\]
The homomorphism is injective in the ordinary outer group,
including all inner automorphisms of the finite quotient.
At \(n=1\), every nonzero parameter gives an actual automorphism
of order five of the characteristic group of order \(5^{19}\).
The parameter maps and their geometric \(\mathfrak S_3\) action
are compatible with coefficient reduction.

These are fixed finite source targets for the full profinite
lifting question.  No assertion that a non-gamma parameter in
\textup{(D.37.2)} has a full profinite peripheral lift is included.
\end{statement}

\begin{proof}
Put \(R_n=\mathbf Z/5^n\), let \(I=(X,Y)\), and define
\[
 B_n=R_n[X,Y]/I^7,\qquad
 U_n=(R_n[X,Y]/I^6)^2.
 \tag{D.37.4}
\]
The ambient finite group consists of pairs \((a,u)\) with
\(a\in1+I B_n\) and \(u\in U_n\), with the exact multiplication
\[
              (a,u)(b,v)=(ab,u+av).
 \tag{D.37.5}
\]
Multiplication by \(a\) on \(U_n\) uses its reduction modulo
\(I^6\).  This is a finite five-group.
Define \(P_n\) to be the actual subgroup generated by
\[
 x=(1+X,(1,0)),\qquad
 y=(1+Y,(0,1)).
 \tag{D.37.6}
\]
Thus \(\rho_n\) is an actual continuous epimorphism from
the free profinite group, rather than a presentation whose
realization is left implicit.

Every element \((a,(u_1,u_2))\) in this source image satisfies
\[
                         a-1=Xu_1+Yu_2
                         \quad\text{in }B_n.
 \tag{D.37.7}
\]
The identity holds for the two generators and is preserved
by \textup{(D.37.5)} and inversion.
The two diagonal generators have order \(5^{n+1}\) and
are independent.  To see the order exactly, for \(k\geq1\)
the coefficients of
\((1+X)^{5^k}-1\) through degree six have valuations
\(k-v_5(j)\) at degree \(j\), when that coefficient is
present.  The degree-five coefficient is nonzero modulo
\(5^n\) for \(k=n\), and all coefficients vanish for
\(k=n+1\).  Specializing respectively \(Y=0\) and
\(X=0\) proves independence.
The same binomial calculation applied to
\(\sum_{r=0}^{5^{n+1}-1}(1+X)^r\), through degree five,
shows that the full group element \(x\), and likewise
\(y\), has order \(5^{n+1}\).

Write \(c=[x,y]\).  Direct multiplication gives
\[
            c=(1,(-Y,X)),\qquad
 c^p:=(1,p(-Y,X))
       \quad\bigl(p\in R_n[X,Y]_{\leq4}\bigr).
 \tag{D.37.8}
\]
Here the notation \(c^p\) denotes the displayed module
element, not an undefined exponent by a polynomial.
It is realized by actual group operations: conjugation
by \(x\) multiplies its coefficient by \(1+X\), and
\([x,c^p]=c^{Xp}\); similarly \([y,c^p]=c^{Yp}\).
Integer powers and products realize every coefficient
in \(R_n\).  Thus all fifteen monomials of total degree
at most four in \textup{(D.37.8)} belong to the derived
subgroup.

Conversely, if the diagonal of a source element is one,
\textup{(D.37.7)} is the actual polynomial identity
\(Xu_1+Yu_2=0\), of degree at most six.
It implies \(u_1=-Yp,\ u_2=Xp\) for a unique polynomial
of degree at most four.  This uses only that multiplication
by a polynomial variable is injective in a polynomial
ring over \(R_n\): setting \(X=0\) first shows that
\(X\) divides \(u_2\), and the identity then gives the
claim.  No division by a nonunit of \(R_n\) occurs.
Therefore the kernel of the diagonal map is precisely
the fifteen-coordinate module in \textup{(D.37.8)}.
It is the derived subgroup.  Every source element has
a unique normal form
\[
 x^r y^s c^p,\qquad
 0\leq r,s<5^{n+1},\quad
 p\in R_n[X,Y]_{\leq4}.
 \tag{D.37.9}
\]
The product of the diagonal and module orders proves
\textup{(D.37.1)} and the asserted abelianization.
Furthermore
\[
 \gamma_jP_n
     =\{c^p:p\in I^{j-2}/I^5\}
       \quad(2\leq j\leq6),\qquad
 \gamma_7P_n=1.
 \tag{D.37.10}
\]
Both inclusions follow from the exact commutators
\([x,c^p]=c^{Xp}\), \([y,c^p]=c^{Yp}\).
In particular the class is exactly six.

\paragraph{The actual fully invariant kernel map.}
Let \(g=(a_g,u_g)\) and \(h=(a_h,u_h)\) be any two elements
of \(P_n\).  The substitutions
\[
 \sigma_{g,h}(X)=a_g-1,\qquad
 \sigma_{g,h}(Y)=a_h-1
 \tag{D.37.11}
\]
preserve the augmentation ideal and its powers.
Define
\[
 \Phi_{g,h}(a,(u_1,u_2))
    =\left(\sigma_{g,h}(a),\
      \sigma_{g,h}(u_1)u_g+
      \sigma_{g,h}(u_2)u_h\right).
 \tag{D.37.12}
\]
This is a homomorphism for
\textup{(D.37.5)}: the module formula is
\(\sigma(u+av)\) followed by the displayed linear
map, and it equals the module part of the product
of the two images.  The truncations have the
required types because \(\sigma(I^j)\subset I^j\).
The map sends \(x\) to \(g\) and \(y\) to \(h\).
Consequently, for every actual word \(w\),
\[
                  w(g,h)
                      =\Phi_{g,h}(\rho_n(w)).
 \tag{D.37.13}
\]
The same equality holds for profinite words by
continuity into the finite groups.
If \(\rho_n(w)=1\), its image under every endomorphism
of \(\widehat F\) is therefore again in
\(\ker\rho_n\).
This proves full invariance, and hence characteristicness,
with the exact quotient map.  It also proves that
the geometric permutations of the three peripherals
act on this same target.  Reducing the coefficients
in \textup{(D.37.4)--(D.37.6)} maps the two generators
onto those of \(P_n\), which proves the stated
surjectivity of \(P_{n+1}\to P_n\).

\paragraph{All three special automorphisms.}
For
\[
 Q=AX^2+BXY+CY^2,\quad
 p=(X+Y)Q,\quad
 f_Q=c^p,\quad h_Q=c^{YQ},\quad
 \Delta_Q=c^{-Y(X+Y)Q},
 \tag{D.37.14}
\]
all four elements on the right are actual group
elements by \textup{(D.37.8)}.
The last belongs to \(\gamma_6P_n\) and is central.
The linear coefficient of \(Y\) in a diagonal
element is an additive homomorphism
\[
 s:P_n\longrightarrow R_n,\qquad s(x)=0,\quad s(y)=1.
 \tag{D.37.15}
\]
Define the actual automorphism
\[
                         \beta_Q(g)=g\Delta_Q^{\,s(g)}.
 \tag{D.37.16}
\]
Centrality proves the homomorphism identity.
It fixes \(\Delta_Q\), since \(s(\Delta_Q)=0\);
the inverse is obtained by replacing \(Q\) by \(-Q\).
The same formula proves
\(\beta_Q\beta_{Q'}=\beta_{Q+Q'}\).
It fixes the whole derived subgroup and fixes \(x\).
The exact relation
\[
           [f_Q,y]=c^{-Yp}=\Delta_Q
 \tag{D.37.17}
\]
shows that \(\beta_Q(y)=f_Qyf_Q^{-1}\).

Put \(z=(xy)^{-1}\).  Then
\(\beta_Q(z)=z\Delta_Q^{-1}\).
Conjugation of a module element by \(z\) multiplies
its coefficient by \(((1+X)(1+Y))^{-1}\).
The identity
\[
 \left(1-((1+X)(1+Y))^{-1}\right)YQ
                =Y(X+Y)Q\pmod{I^5}
 \tag{D.37.18}
\]
is exact at the required degree: the difference
between the first parenthesis and \(X+Y\) has
degree at least two, whereas \(YQ\) has degree three.
It follows that
\([h_Q,z]=\Delta_Q^{-1}\), and hence
\[
                         \beta_Q(z)=h_Qzh_Q^{-1}.
 \tag{D.37.19}
\]
Thus all three peripheral conditions hold as
literal identities in \(P_n\).
The image of the normalized conjugator in
\(\gamma_5P_n/\gamma_6P_n\) is precisely
\((X+Y)Q[x,y]\), because the group conjugation
differences \(X,Y\) induce the actual ordinary
graded brackets.  This proves the stated leading
coordinate comparison.
The polynomial \(-Y(X+Y)Q\) has degree four;
multiplication by \(Y(X+Y)\) is injective on
the quadratic coefficient module.  Thus
the order of \(\beta_Q\) is exactly the additive
order of \(Q\), before taking outer classes.

\paragraph{The ordinary outer quotient, including restricted inner powers.}
We compute the actual centralizer of \(x\), rather
than using a centralizer theorem for a different
free group.  Let
\(g=x^r y^s c^p\) be the normal form
\textup{(D.37.9)}, and write
\[
                   S_s(Y)=\sum_{j=0}^{s-1}(1+Y)^j.
 \tag{D.37.20}
\]
The second module coordinate of \(g\) is
\((1+X)^r(S_s(Y)+(1+Y)^sXp)\).
Commuting with \(x\) requires its product by
\(X\) to vanish modulo \(I^6\).
Setting \(X=0\) in the coefficient condition
shows that all coefficients of \(S_s(Y)\)
through degree four vanish modulo \(5^n\).
The constant coefficient first forces
\(5^n\mid s\).  If \(s=5^n t\), \(0\leq t<5\),
its degree-four coefficient is
\[
                    \binom{5^n t}{5}
                       \equiv5^{n-1}t\pmod{5^n}.
 \tag{D.37.21}
\]
Hence \(s=0\).
The remaining second-coordinate equation is
\(X^2p=0\pmod{I^6}\), which forces \(p\)
to be homogeneous of degree four.
Conversely such \(c^p\) is central.
This proves the exact equality
\[
                        C_{P_n}(x)=
                        \langle x\rangle\gamma_6P_n.
 \tag{D.37.22}
\]
Intersecting with the centralizer of \(y\)
also proves \(Z(P_n)=\gamma_6P_n\).

Suppose \(\beta_Q\) is an inner automorphism.
Since it fixes \(x\), \textup{(D.37.22)}
reduces that inner automorphism to conjugation
by \(x^r\).  Its effect on \(y\) is
\[
                         [x^r,y]=c^{S_r(X)}.
 \tag{D.37.23}
\]
For this to have degree four, the lower
coefficients force \(r=5^n t\); its coefficient
is then \(5^{n-1}tX^4\) modulo \(I^5\).
On the other hand
\(\Delta_Q=c^{-Y(X+Y)Q}\) has coefficient
divisible by \(Y\).
Equality forces \(t=0\), followed by \(Q=0\).
This accounts explicitly for the inner powers
of \(x\) that are trivial in lower degrees
but can act in degree five.
It proves the ordinary outer injectivity
in \textup{(D.37.2)}.

\paragraph{Symmetry and executable finite comparisons.}
The actual source substitutions
\[
 (x,y)\longmapsto(y,x),\qquad
 (x,y)\longmapsto(y,(xy)^{-1})
 \tag{D.37.24}
\]
act on \(P_n\) by \textup{(D.37.12)} and have
orders two and three.  Conjugate \(\beta_Q\)
by either substitution.  Its displacement
of \(x\) is \(c^{r(X,Y)}\), with \(r\)
homogeneous of degree four and divisible
by \(X\).  Conjugating further by the actual
element \(c^{r/X}\) removes this displacement.
The remaining automorphism is
\[
 \beta_{Q(Y,X)}\quad\text{or}\quad
 \beta_{Q(Y,-X-Y)},
 \tag{D.37.25}
\]
respectively.
These identities follow directly from
\textup{(D.37.14)--(D.37.18)}; every
displacement lies in the central sixth
layer, so terms of higher degree in the
source substitutions vanish.
They are the exact finite-group comparison
with the symmetric-square action in D.21,
with its inner correction retained.
The line \(Q=X^2+XY+Y^2\) is the normalized
gamma polynomial line.  This identifies
the polynomial line; it does not assert
that its primitive finite parameter is
realized by a Galois element.

The program
\texttt{a\_weight5\_characteristic.py}
and its saved output and JSON certificate
execute the constructions for \(n=1,2\).
They check all 625 and 15625 distinct
diagonal elements, all fifteen actual
commutator-module words, all three
special maps and their third peripheral
conjugators, and the complete 125-parameter
injection at \(n=1\).
The finite centralizer exponent conditions
and both geometric symmetry corrections
are checked as well.
The resulting orders are
\[
 |P_1|=19073486328125,\qquad
 |P_2|=14551915228366851806640625.
 \tag{D.37.26}
\]
The proof above supplies every precision
\(n\); it does not extrapolate a group
or kernel identity from these two runs.
In particular the maps
\(\operatorname{Out}^{\sharp}(\widehat F)
\to\operatorname{Out}(P_n)\) now have
fully specified characteristic targets
and fully specified order-five test
elements at \(n=1\).
Deciding which compatible parameters
lift through all finer finite quotients
remains the full-source bridge.
\end{proof}
